\clearpage
\phantomsection

\addcontentsline{toc}{chapter}{Tóm tắt}
\chapter*{\fontsize{13}{13}\selectfont\textbf{TÓM TẮT}}

\noindent\textbf{Tóm tắt:} Khóa luận nghiên cứu và triển khai một tác nhân học tăng cường sâu cho game hành động Roguelike góc nhìn từ trên xuống trong Unity. Tác nhân chính, CombatAgent, sử dụng quan sát vector 207 chiều và không gian hành động rời rạc nhiều nhánh $(3,3,3,9)$ để phối hợp di chuyển, chiến đấu và hướng tấn công. Môi trường huấn luyện được mở rộng bằng quy trình PCGNN sinh bản đồ ngoại tuyến, bộ chuyển bản đồ sang định dạng của game, bộ kiểm tra khả năng đi tới và cơ chế chia bản đồ thành các mức Dễ, Trung bình và Khó phục vụ học theo lộ trình (\emph{curriculum learning}, CL).

Trên nền PPO của Unity ML-Agents, khóa luận thực hiện ma trận thí nghiệm cô lập biến gồm chín cấu hình: ICM, RND, LSTM, CL và biến thể hàm thưởng. Trong các run đã thực hiện, CL cải thiện rõ nhất: phần thưởng tăng từ 2.073 ở cấu hình PPO cơ sở lên 3.843 ở vùng Khó (+85\%). Bổ sung tín hiệu dẫn đường BFS vào quan sát (run59) làm tỷ lệ hoàn thành phòng tăng mạnh, cho thấy điều hướng là nút thắt chính của môi trường nhưng không chứng minh tác nhân tự học tìm đường hoàn toàn. Cấu hình mở rộng run62 -- quan sát lưới không gian cục bộ, CL 4 bậc, ICM và huấn luyện dài 40.82M bước -- đạt tỷ lệ dọn phòng 0.666 và entropy 3.048, vượt qua vùng bão hòa entropy 4.40 nat; do thay đổi nhiều yếu tố cùng lúc nên kết quả này chỉ mang tính thăm dò.

Do giới hạn tài nguyên, mỗi cấu hình mới chỉ được chạy một seed và quy mô huấn luyện giữa nhóm không dùng CL và nhóm dùng CL chưa hoàn toàn đồng nhất. Vì vậy, các kết luận định lượng được diễn giải như xu hướng thực nghiệm ban đầu, đồng thời khóa luận đề xuất các hướng mở rộng như chạy nhiều seed, bổ sung mốc so sánh dựa trên luật, đánh giá trên bản đồ chưa từng thấy và thiết kế lại không gian hành động.

\vspace{0.5cm}
\noindent{\textbf{Từ khóa:}} \textit{Học tăng cường sâu, Game Roguelike, Procedural Content Generation (PCG), PPO (Proximal Policy Optimization), học theo lộ trình (CL), tác nhân tự chủ}
